\section{Data for nonlinear and time dependent Options} \label{sec2.13}

\textbf{General remarks}

\medskip %

The data in this block are used to define a discretization in
time, and the test-particle self interaction effects. These latter
options are based upon the ``Bird's \ac{DMCS}"
procedure, and are currently not in use (see
\cite{kn:Behringer} for a detailed description of its
implementation in EIRENE). They can be activated by replacing the
current dummy routine STOSS, which is called after each time-step
from subroutine EIRENE, by a routine that reads the particle
population from the ``census arrays" described below, and which
then carries out the binary self-collisions (modification of the
individual particle velocity vectors). The rest of the code for
time dependency and the \ac{DMCS} algorithm is the same and described
in this section. Note that, currently, with the use of a dummy
routine STOSS, non-linear self collision effects can still be
simulated, in \ac{BGK} approximation, by using the iterative mode of
operation (NITERI, NITERE, input block 1), see
\cref{sec1.8}.

As for the time-dependence options, it is important to note that
there are two kinds of time-steps. One is the so called
``time-cycle". Each time cycle is a more or less complete EIRENE
run in its own, even if several (NTIME, see \cref{sec2.1})
such time-cycles are carried out in one single job. In each
time-cycle, a time-horizon is set and the trajectories are stopped
(at latest) if they have survived until then. All relevant
particle coordinates (position, velocity, type and species, cell
indices, etc.) are stored at this time-horizon on the so called
``census-arrays".

Each time cycle may consist of one or more than one (NTMSTP, see
below) ``internal time steps". At the end of each of these
time-steps, the relevant snapshot tallies are updated (subroutine TIMCOL).
These are the volumetric ``snapshot tallies", the default ``time-surface
tallies" (particle and energy fluxes) and the ``census arrays". Hence: these tallies are averaged
over NTMSTP (internal) time-steps. For more details on snapshot
tallies see \cref{sec3.2.3}.

This distinction between complete ``time cycles" and ``internal
time steps" is particularly relevant for obtaining truly steady
state results on the census arrays, e.g., in order to continue a
stationary case in the time-dependent mode (see below).

 \medskip

\textbf{The Input Block}

\begin{lstlisting}
*** 13. Data for nonlinear and time dependent Options
      NPRNLI NINITL_READ NPRMUL
      IF (NLERG.AND.NPRNLI.EQ.0) NPRNLI=100
      IF (NPRNLI.GT.0) THEN
        NPTST,NTMSTP
        DTIMV,TIME0
** 13A. DATA FOR SNAPSHOT TALLIES
        NSNVI
        IF (NSNVI.GT.NSNV)
          DO 1320 J=1,NSNVI
            ISNVE(J),ISNVS(J),ISNVT(J),ISNRC(J)
            TXTTAL(J,NTALT)
            TXTSPC(J,NTALT),TXTUNT(J,NTALT)
 1320     CONTINUE
        ENDIF
\end{lstlisting}

\medskip

\textbf{Meaning of the Input Variables for this block}

\medskip

\begin{description}
\item[NPRNLI] Total number of test particles in time
dependent arrays (``census arrays") (in old versions before 2001: NPRNLI must be $\le$ NPRNL,
see: PARMUSR). The scoring on census arrays stops at latest when NPRNLI scores are on
the census array.

If NPRNLI $>0$, but the rest of the data in this block are not
specified, then a default time-horizon is defined. See default values
specified below.

If NPRNLI $=0$, no census arrays are scored

In case NLERG = TRUE  (see input block 1), a time horizon is
absolutely necessary in order to prevent infinite histories.
Therefore, in case NLERG=TRUE and NPRNLI=0 an automatic correction
to NPRNLI=100 is carried out.
\item[NINITL\_READ]  (new: 2013) Same as NINITL in block 7: provides random number seed for ``time-stratum" (sampling from census
array  (default: =0, no fresh initialization of random number generator for this stratum)
\item[NPRMUL] (new: 2013) Multiplicative factor for NPRNLI, in order to increase size of census
to more than 999999 particles, which otherwise would be the maximum due to I6 formatted integer input. 
Default: = 0: no multiplication carried out
\item[NPTST] Same as NPTS in block 7.
This is the number of histories, which are continued from a
previous time-cycle (``Time dependence stratum ISTRA=NSTRAI+1").
The initial coordinates are randomly sampled (with replacement)
from the census-array data from an earlier time-cycle. The
probability for sampling a particular particle from the census
array is proportional to its weight stored on the census array as
well. Due to ``sampling with replacement" an individual particle,
which is on census, may be sampled more than once, or not at all,
with the likelihood for these events given by its weight (``warm
restart"). This census array is either defined at the end of the
previous time cycle in the same run (subr.\ TMSTEP), or it is read
from an earlier run from stream 15 (via a call to subr.\ RSNAP from
subr.\ INPUT) in the initial phase of the run, for the very first
time-cycle (continuation of an earlier sequence of time-cycles).

If NPTST $=$ 0, then NPTST is reset to IPRNL. IPRNL is the
the number of scores on the census array in the previous time cycle.

If NPTST $<$ 0, then NPTST is reset to IPRNL, and the random
sampling from the census array is now replaced by a one-to-one
re-launch of all particles from the census array without random
sampling (``cold restart"). Until Aug.\ 2015 this option was available only in connection with the
NLMOVIE option (movies of trajectories)
and had led to other modifications of the run parameters as well.
(Automatically then internally: NLMOVIE = TRUE). New: NLMOVIE AND NPTST $<$ 0 options
are now independent from each other. In both cases: one by one re-launch from old census
is enforced, rather than random sampling from old census.


Default: NPTST=0
\item[NTMSTP]
Total number of time-steps for particle tracing.
Each trajectory can score on census up to NTMSTP times.
Particle trajectories are stopped after NTMSTP time-steps.

Default: NTMSTP=1

For convenience and by abuse of language, we refer to the
3-dimensional hyper-surface $t=t_n$ of the four dimensional
$(\vr,t)$-space as ``time-surface", and, hence, tallies evaluated
at fixed time $t_n$ (``snapshot-tallies") are surface averaged
tallies in this terminology.

Fluxes onto this surface are stored on the arrays for a surface
no.\ NLIM+NSTSI+1, which is added automatically to the NLIM
additional and NSTSI non-default standard surfaces.

The time-surface is transparent for NTMSTP-1 steps
and absorbing afterwards, i.e., absorbing at time t = NTMSTP $\times$ DTIMV.
Snapshot tallies are averages over NTMSTP time-steps.

Default: NTMSTP=1

If NTMSTP $<$ 0, then each particle can score an unlimited number of
times on census, i.e., the time-surface is always transparent. This
option can be used for initialization of census arrays for time
dependent runs. The census arrays then represent a stationary
distribution corresponding to a certain constant (in time) influx of
particles, rather than an estimate at a fixed time. Hence, for any
fixed detector function (volume averaged tally), the snapshot
estimator should give the same results (up to statistical precision)
as the track-length estimator or the collision estimator. Note that
in order to obtain this stationary estimate from snapshot estimates
an additional multiplicative factor DTIMV (\si{\second}) (see next) is applied
to snapshot tallies in case NTMSTP $<$ 0.
\item[DTIMV]
Length of each individual internal time-step (\si{\second})

Default: DTIMV=1.D-02

\item[TIME0]
Initial time $t_0$ of the first time-step.
(irrelevant, only for printout and book-keeping)

Default: TIME0=0.

\item[NSNVI]
Number of snapshot tallies computed from census arrays. In old
EIRENE versions without dynamic allocation of storage NSNVI must be
less or equal NSNV (see PARMUSR), and the detector functions are
user supplied in subroutine UPNUSR, see Sub-\cref{sec3.2.3}.

Default: NSNVI=0

\end{description}
The meaning of the next three cards is the same as for the corresponding
cards in block 10A, 10B or 10D.

\bigskip
\textbf{Census Arrays}

All particle co-ordinates at time $t_i$ for history number IHIST,
i.e., position, cell indices, velocity etc.\ are stored on arrays
RPART(IHIST, \dots) (Real) and IPART(IHIST, \dots) (Integer) in
subroutine TIMCOL. This subroutine TIMCOL is called if ``a
collision with the time surface $t=t_i$" has occurred.

The number of scores on the old census arrays (from a previous completed time-cycle
in a current run or from an earlier run) is \textbf{IPRNL}.

The actual number of a score in subroutine TIMCOL is \textbf{IPRNLI},
and after the score IPRNLI is increased by one. IPRNLI is set to
zero at the beginning of a run.

The actual number of a score in subroutine TIMCOL for the present
stratum ISTRA  is \textbf{IPRNLS}, and after
the score IPRNLS is increased by one. IPRNLS is reset to zero at the
beginning of sampling for each stratum.

Scaling of census array fluxes:   to be written
